\documentclass[10pt,a4paper]{amsart}
\usepackage[margin=19mm]{geometry}
\usepackage{amsmath,amssymb,mathtools}
\usepackage[scr=rsfs,cal=euler]{mathalfa}
\usepackage{xfrac}
\usepackage[unicode,hidelinks,bookmarks=false]{hyperref}
\newcommand{\ie}{i.e.\ }
\newcommand{\cf}{cf.\ }
\newcommand{\resp}{respectively\ }
\newcommand{\Subsection}{\subsection}
\providecommand{\nobreakdash}{\nobreak\hbox{-}}
\setlength{\emergencystretch}{2em}
\multlinegap0pt
\usepackage{latexsym,amscd,mathrsfs,url}
\usepackage{color}
 \newcommand{\sgn}{\operatorname{sgn}}
\newcounter{count}
\numberwithin{count}{section}
\newtheorem{Lemma}[count]{Lemma}
\newtheorem{Remark}[count]{Remark}
\newtheorem{Example}[count]{Example}
\newtheorem{Definition}[count]{Definition}
\newtheorem{theorem}[count]{Theorem}
\newtheorem{Corollary}[count]{Corollary}
\newtheorem{Theorem}[count]{Theorem}
\newtheorem{Conjecture}[count]{Conjecture}
\newtheorem{Problem}[count]{Open Problem}
\newtheorem{Statement}[count]{Statement}
\begin{document}
\title[Polynomials interpolated totally positive sequences]{Polynomials interpolated totally positive sequences}
\author{Anna Vishnyakova}
\date{}
\maketitle
\noindent\emph{Source and licence.} Polynomials interpolated totally positive sequences. Version arXiv:2607.06207v1 (CC BY 4.0). This material is distributed under the Creative Commons Attribution 4.0 International licence, http://creativecommons.org/licenses/by/4.0/, as recorded in the arXiv OAI licence metadata for this exact version and shown on the abstract page https://arxiv.org/abs/2607.06207v1. Full rights notice: licenses/arxiv-2607.06207-CC-BY-4.0.txt.\par\smallskip
\noindent\textit{Editorial scope.} Counted unit: the original \texttt{Theorem} environment \texttt{th1} at source lines 373--399 together with its complete proof at lines 401--468; the delivered document keeps source lines 317--469 so that every referenced label and every citation resolves inside it. Native editable source is the paper's own concatenated TeX; the AMS wrapper is editorial formatting only. No publisher class, upstream script or unrelated artwork is redistributed.
\begin{equation}
\label{fd2}
\mathcal{L}(P(-x-1)) (x) = Q(-x-1).
\end{equation}
Note that if $Q$ has only real zeros located in the segment $[-1, 0],$ 
then $Q(-x-1)$ also  has only real zeros located in the segment $[-1, 0].$
Thus, we have proved the following statement.
 
\begin{Statement}
\label{St0}
Let $P$ be a real polynomial of degree $n\in \mathbb{N}.$  If $(P(k))_{k=0}^\infty \in \mathrm{TP},$ then $\left((-1)^n 
P(-k-1)\right)_{k=0}^\infty \in \mathrm{TP}.$
\end{Statement}

\section{Necessary conditions}

Suppose that a polynomial 
$$P(x)= a_0 + a_1 x +a_2 \frac{x(x-1)}{2!} 
+  \ldots +a_n \frac{x(x-1)\cdot \ldots \cdot (x-n+1)}{n!}$$
is such that
$$\mathcal{L}(P)
= a_0 +a_1 x +a_2 x^2 + \ldots + a_n x^n$$ 
has only real zeros located in the segment $[-1, 0].$ We have mentioned  the 
necessary condition that all  coefficients of $\mathcal{L}(P)$ are of the same signs. 

The following statement gives some more simple necessary conditions.

\begin{Statement}
\label{st1}
Suppose that a real polynomial 
$$P(x)= a_0 + a_1 x +a_2 \frac{x(x-1)}{2!} 
+  \ldots +a_n \frac{x(x-1)\cdot \ldots \cdot (x-n+1)}{n!},$$
$a_n=1, $   is such that $(P(k))_{k=0}^\infty \in \mathrm{TP}.$ Then

1. All real zeros of $P$ belong to the segment $[- n, n-1].$

2. For every $j=0, 1, 2, \ldots, n-1$ we have $a_j \leq C_n^j.$

3. For every $j= 1, 2, \ldots, n-1$ such that $a_{j-1}>0$  we have
$$\frac{a_j^2}{a_{j-1}a_{j+1}} \geq \frac{j+1}{j} \cdot \frac{n-j+1}{n-j}$$
(Newton inequalities).

\end{Statement}

\begin{proof}

The statement that all non-negative zeros of $P$ are in the segment $[0, n-1]$
follows from the fact that all coefficients of $P$ are of the same signs.
The statement that all non-positive zeros of $P$ are in the segment $[-n, 0]$
follows from the previous statement and formula~(\ref{fd2}). Statements~2 and 3
are well-known necessary conditions for a polynomial to have all real zeros
in the statement $[-1, 0].$
\end{proof}

It is interesting to compare the following necessary conditions with Theorem~B.


\begin{Theorem}
\label{th1}
Suppose that a real polynomial 
$$P(x)= a_0 + a_1 x +a_2 \frac{x(x-1)}{2!} 
+  \ldots +a_n \frac{x(x-1)\cdot \ldots \cdot (x-n+1)}{n!},$$
$a_n=1, $   is such that $(P(k))_{k=0}^\infty \in \mathrm{TP}.$ 

1. If there exists $k=1, 2, \ldots, n-1$ such that $P(k)=0,$ then
$P(0) = P(1) = \ldots = P(k-1)=0.$

2.  For every $m\in \mathbb{N}$ we get
\begin{equation}
\label{f110}
P(- m) = \frac{1}{(m-1)!}  \left(x^{m-1} Q(x)\right)^{(m-1)}|_{x=-1}.
\end{equation}

3.  If there exists $l \in \mathbb{N}$ such that $P(-l)=0,$ then
$P(-1) = P(-2) = \ldots = P(-(l-1))=0.$

4. If $\deg P = 2l+1, l \in \mathbb{N}\cup \{0\},$  then there is a root of
$P$ in the segment $[-1, 0].$

%4. Let $P \in \mathbb{R}[x]$  be a polynomial with positive leading coefficient 
%and $x_1 \leq x_2 \leq \ldots \leq x_m$ be all real zeros of $P.$  Suppose that  
%$(P(k))_{k=0}^\infty \in \mathrm{TP}.$ Then for every $j=1, 2, \ldots, m-1$ we have 
%$x_{j+1} - x_j \leq 1.$   
\end{Theorem}

\begin{proof}

1. We have for all $k=1, 2, \ldots, n-1$
\begin{equation}
\label{f7}
P(k) = \sum_{j=0}^k C_k^j  a_j.
\end{equation}
If k $P(k)=0,$ then, since all coefficients $a_j$ are of the same signs, we obtain
$a_0=a_1= \ldots = a_k =0$, thus $P(0) = P(1) = \ldots = P(k-1)=0.$



2. For  $P(x)= a_0 + a_1 x +a_2 \frac{x(x-1)}{2!} 
+  \ldots +a_n \frac{x(x-1)\cdot \ldots \cdot (x-n+1)}{n!},$ 
we denote by  $Q(x) = \mathcal{L}(P) (x) = a_0 +a_1x +a_2 x^2 + \ldots + a_n x^n.$
We have
\begin{equation}
\label{f8a}
P(0) = a_0  = Q(0).
\end{equation}
\begin{equation}
\label{f8}
P(-1) = a_0 - a_1 + a_2 -a_3 + \ldots + (-1)^n a_n = Q(-1).
\end{equation}
\begin{equation}
\label{f9}
P(-2) = a_0 - 2 a_1 + 3a_2 - \ldots + (-1)^n (n+1) a_n = (x Q(x))^\prime|_{x=-1}.
\end{equation}
Since $(P(k))_{k=0}^\infty \in \mathrm{TP},$ all zeros of $Q$ are real and in the segment $[-1, 0].$
Thus, all zeros of the polynomial $xQ(x)$ are real and in the segment $[-1, 0].$ Suppose 
that $P(-2)=0.$  That means that  $(x Q(x))^\prime|_{x=-1} =0,$  whence  $-1$ is a root of $Q$ 
of multiplicity $\geq 2.$ In particular, $Q(-1) = P(-1) =0.$

For every $m\in \mathbb{N}$ we get
$$
P(- m) = \frac{1}{(m-1)!} \left( \frac{(m-1)! }{0!} a_0 - \frac{m!}{1!} a_1 + \frac{(m+1)!}{2!} a_2 -  \frac{(m+2)!}{3! } a_3  \right.$$   
$$\left.   + \ldots + (-1)^n \frac{(m+n-1)!}{n!} a_n\right ) =
\frac{1}{(m-1)!}  \left(x^{m-1} Q(x)\right)^{(m-1)}|_{x=-1}.
$$

3.  Since $(P(k))_{k=0}^\infty \in \mathrm{TP},$ all zeros of $Q$ are real and in the segment $[-1, 0].$
Thus, all zeros of the polynomial $x^{l-1} Q(x)$ are real and in the segment $[-1, 0].$ Suppose 
that $P(- l)=0.$  That means that  $\left(x^{l-1} Q(x)\right)^{(l-1)}|_{x=-1} =0,$  whence  $-1$ 
is a root of $Q$ of multiplicity $\geq l.$ Using (\ref{f110}), we obtain that  
$P(-1) = P(-2) = \ldots = P(-(l-1))=0.$


4. Let  $\deg P = \deg Q = 2l+1, l \in \mathbb{N}\cup \{0\}. $ Since $(P(k))_{k=0}^\infty \in \mathrm{TP},$ 
all zeros of $Q$  are real and in the segment $[-1, 0].$  Thus, $Q(0)\cdot Q(-1) \leq 0.$
By (\ref{f8a}) and (\ref{f8}) we get $P(0)\cdot P(-1) \leq 0.$ Hence, then there is a 
root of $P$ in the segment $[-1, 0].$

%4. Suppose that  there exists $j=1, 2, \ldots, m-1$ such that
%$x_{j+1} - x_j  > 1.$   Without loss of generality we can assume that
%$x_{j+1} >0$  (otherwise we consider $P(-x-1)$ instead of $P(x),$ and use 
%the formula (\ref{fd2})). Since $x_{j+1} - x_j  > 1,$ there exists $k=0, 1, 2, \ldots$
%such that $x_{j} < k < x_{j+1}.$ By (\ref{f7}) we have $P(k) = \sum_{j=0}^k C_k^j  
%a_j \geq 0, $  and, since $x_{j} < k < x_{j+1},$ we get $P(k) >0.$
%For  $P(x)= a_0 + a_1 x +a_2 \frac{x(x-1)}{2!} 
%+  \ldots +a_n \frac{x(x-1)\cdot \ldots \cdot (x-n+1)}{n!},$ we have
%\begin{equation}
%\label{diff}
%P(x+1) - P(x)= a_1 + a_2 x +a_3 \frac{x(x-1)}{2!} 
%+  \ldots +a_n \frac{x(x-1)\cdot \ldots \cdot (x-n+2)}{(n-1)!}.
%\end{equation}


\end{proof}


\end{document}
